\documentclass[11pt]{article}

\usepackage[a4paper,margin=1in]{geometry}
\usepackage{amsmath,amssymb}
\usepackage{graphicx}
\usepackage{booktabs}
\usepackage{array}
\usepackage{xcolor}
\usepackage{listings}
\usepackage{hyperref}
\usepackage{enumitem}
\usepackage{fancyhdr}
\usepackage{tcolorbox}
\usepackage{caption}
\usepackage{tabularx}

\hypersetup{colorlinks=true, linkcolor=blue, urlcolor=blue, citecolor=blue}

\lstdefinestyle{pystyle}{
    language=Python,
    basicstyle=\ttfamily\footnotesize,
    keywordstyle=\color{blue},
    commentstyle=\color{gray},
    stringstyle=\color{teal},
    numbers=left,
    numberstyle=\tiny\color{gray},
    stepnumber=1,
    numbersep=8pt,
    backgroundcolor=\color{gray!5},
    frame=single,
    breaklines=true,
    breakatwhitespace=false,
    showstringspaces=false,
    tabsize=4
}
\lstset{style=pystyle}

\pagestyle{fancy}
\fancyhf{}
\setlength{\headheight}{14pt}
\rhead{CS3807 -- Deep Learning Laboratory}
\lhead{Experiment 3}
\rfoot{\thepage}

% Reusable note box for illustrative (not actually re-run) values
\newtcolorbox{estbox}{colback=yellow!8, colframe=orange!70!black, boxrule=0.6pt, arc=1pt,
  left=4pt, right=4pt, top=2pt, bottom=2pt}

\title{\textbf{SHIV NADAR UNIVERSITY CHENNAI} \\[2pt]
\large B.Tech Artificial Intelligence and Data Science \\[2pt]
CS3807 -- DEEP LEARNING LABORATORY \\[10pt]
\Large EXPERIMENT 3 \\[4pt]
\large Implementation of Convolutional Neural Networks (CNNs) for Image Classification}
\author{}
\date{}

\begin{document}
\maketitle
\thispagestyle{fancy}

\begin{center}
\begin{tabular}{|l|l|}
\hline
\textbf{Degree} & B.Tech Artificial Intelligence and Data Science \\
\hline
\textbf{Semester} & V \\
\hline
\textbf{Subject Code} & CS3807 \\
\hline
\textbf{Subject} & Deep Learning Laboratory \\
\hline
\textbf{Academic Year} & 2026--27 \\
\hline
\textbf{Name} & Ashwin S \\
\hline
\textbf{Register Number} & 24011101006 \\
\hline
\textbf{Date} & 12/08/26 \\
\hline
\end{tabular}
\end{center}



\tableofcontents
\newpage

% =========================================================
\section{Objective}
The objective of this experiment is to understand and implement the working principle of Convolutional Neural Networks (CNNs) for image classification using TensorFlow and Keras.

The experiment specifically aims to:
\begin{enumerate}
    \item Understand the convolution operation.
    \item Compute convolution output dimensions using stride and padding.
    \item Study different convolution kernels and their effects on images.
    \item Visualize feature maps generated by convolutional layers.
    \item Implement and compare max pooling and average pooling.
    \item Calculate the number of trainable parameters in convolutional layers.
    \item Construct a CNN for CIFAR-10 image classification.
    \item Train the CNN using the Adam optimizer.
    \item Evaluate the CNN using accuracy, precision, recall, F1-score, confusion matrix and classification report.
    \item Analyze the effects of kernel size, pooling strategy and number of filters on CNN performance.
\end{enumerate}

% =========================================================
\section{Learning Outcomes}
After completing this experiment, the student will be able to:
\begin{enumerate}
    \item Explain the working principle of convolutional neural networks.
    \item Perform convolution operations using numerical matrices.
    \item Calculate the spatial dimensions of convolution outputs.
    \item Explain the effect of stride and padding.
    \item Identify the purpose of common image-processing kernels.
    \item Implement max pooling and average pooling.
    \item Calculate trainable parameters in CNN layers.
    \item Visualize intermediate feature maps.
    \item Build and train CNN architectures using TensorFlow/Keras.
    \item Evaluate image classification models using multiple performance metrics.
\end{enumerate}

% =========================================================
\section{Background Theory}
A Convolutional Neural Network (CNN) is a type of deep neural network specifically designed for processing data with a grid-like structure, such as images.

Unlike a fully connected neural network, a CNN exploits the spatial structure of images through local receptive fields and shared weights. This enables the network to learn visual features efficiently while requiring significantly fewer parameters.

A typical CNN consists of the following components:
\begin{itemize}
    \item Convolutional Layer
    \item Activation Function
    \item Pooling Layer
    \item Fully Connected Layer
    \item Output Layer
\end{itemize}

The CNN used in this experiment consists of two convolutional blocks followed by a fully connected classification layer.

\subsection{Convolution Operation}
A convolution operation applies a small matrix called a kernel or filter over the input image. At each spatial location, the kernel values are multiplied element-wise with the corresponding input values and summed to generate an output value.

Mathematically,
\[
Y(i,j) = \sum_m \sum_n X(i+m, j+n)\,K(m,n)
\]
where:
\begin{itemize}
    \item $X$ is the input image or feature map.
    \item $K$ is the convolution kernel.
    \item $Y$ is the resulting feature map.
    \item $i, j$ represent the spatial position of the kernel.
    \item $m, n$ represent positions inside the kernel.
\end{itemize}

\subsection{Output Dimension of Convolution}
The spatial dimension of a convolution output is calculated using:
\[
O = \left\lfloor \frac{N - F + 2P}{S} \right\rfloor + 1
\]
where:
\begin{itemize}
    \item $N$ = input spatial dimension
    \item $F$ = filter/kernel size
    \item $P$ = padding
    \item $S$ = stride
    \item $O$ = output spatial dimension
\end{itemize}

For example, for a $32 \times 32$ input, a $3 \times 3$ kernel, stride 1 and valid padding:
\[
O = \frac{32 - 3 + 0}{1} + 1 = 30
\]
Therefore, the output spatial dimension is: $30 \times 30$.

\subsection{Stride}
Stride specifies the number of pixels by which the convolution kernel moves at each step. A stride of 1 moves the kernel one pixel at a time, while a stride of 2 moves the kernel two pixels at a time.

Increasing the stride reduces the spatial dimensions of the resulting feature map.

\subsection{Padding}
Padding involves adding additional pixels around the boundary of an input image before performing convolution.

Two commonly used padding schemes are:
\begin{itemize}
    \item \textbf{Valid padding:} No padding is applied. The spatial dimensions decrease after convolution.
    \item \textbf{Same padding:} Padding is applied so that the output spatial dimensions are preserved when the stride is 1.
\end{itemize}

\subsection{ReLU Activation}
The Rectified Linear Unit (ReLU) is defined as:
\[
\mathrm{ReLU}(x) = \max(0, x)
\]
ReLU replaces negative values with zero while retaining positive values.

It is widely used in CNN hidden layers because it is computationally simple and generally helps neural networks avoid the severe vanishing-gradient problem associated with saturating activation functions.

\subsection{Pooling}
Pooling reduces the spatial dimensions of feature maps while retaining important information.

Two common pooling methods are:
\begin{itemize}
    \item Max Pooling
    \item Average Pooling
\end{itemize}

Max pooling selects the largest value from each pooling region:
\[
Y(i,j) = \max(X_{\text{window}})
\]
Average pooling computes the mean:
\[
Y(i,j) = \frac{1}{N} \sum_{x \in \text{window}} x
\]

% =========================================================
\section{Common Convolution Kernels}
A convolution kernel is a small matrix used to extract specific visual characteristics from an image. Different kernels emphasize different properties such as edges, textures, smooth regions and fine details.

\subsection{Identity Kernel}
\[
K = \begin{bmatrix} 0 & 0 & 0 \\ 0 & 1 & 0 \\ 0 & 0 & 0 \end{bmatrix}
\]
The identity kernel preserves the original image and is useful as a baseline for understanding convolution.

\subsection{Sobel-X Kernel}
\[
K = \begin{bmatrix} -1 & 0 & 1 \\ -2 & 0 & 2 \\ -1 & 0 & 1 \end{bmatrix}
\]
The Sobel-X kernel detects vertical edges by measuring changes in intensity along the horizontal direction.

\subsection{Sobel-Y Kernel}
\[
K = \begin{bmatrix} -1 & -2 & -1 \\ 0 & 0 & 0 \\ 1 & 2 & 1 \end{bmatrix}
\]
The Sobel-Y kernel detects horizontal edges.

\subsection{Laplacian Kernel}
\[
K = \begin{bmatrix} 0 & -1 & 0 \\ -1 & 4 & -1 \\ 0 & -1 & 0 \end{bmatrix}
\]
The Laplacian kernel detects edges in multiple directions.

\subsection{Sharpening Kernel}
\[
K = \begin{bmatrix} 0 & -1 & 0 \\ -1 & 5 & -1 \\ 0 & -1 & 0 \end{bmatrix}
\]
The sharpening kernel emphasizes high-frequency image details.

\subsection{Box Blur Kernel}
\[
K = \frac{1}{9}\begin{bmatrix} 1 & 1 & 1 \\ 1 & 1 & 1 \\ 1 & 1 & 1 \end{bmatrix}
\]
The box blur kernel performs image smoothing by averaging neighboring pixels.

\subsection{Gaussian Blur Kernel}
\[
K = \frac{1}{16}\begin{bmatrix} 1 & 2 & 1 \\ 2 & 4 & 2 \\ 1 & 2 & 1 \end{bmatrix}
\]
Gaussian blur reduces noise while generally preserving image structure better than a simple uniform averaging filter.

\subsection{Emboss Kernel}
\[
K = \begin{bmatrix} -2 & -1 & 0 \\ -1 & 1 & 1 \\ 0 & 1 & 2 \end{bmatrix}
\]
The emboss kernel emphasizes directional intensity changes to create a three-dimensional visual effect.

\subsection{Outline Detection Kernel}
\[
K = \begin{bmatrix} -1 & -1 & -1 \\ -1 & 8 & -1 \\ -1 & -1 & -1 \end{bmatrix}
\]
The outline kernel emphasizes object boundaries and regions with rapid intensity changes.

\subsection{Motion Blur Kernel}
\[
K = \frac{1}{5}\begin{bmatrix}
1 & 0 & 0 & 0 & 0 \\
0 & 1 & 0 & 0 & 0 \\
0 & 0 & 1 & 0 & 0 \\
0 & 0 & 0 & 1 & 0 \\
0 & 0 & 0 & 0 & 1
\end{bmatrix}
\]
The motion blur kernel simulates diagonal motion blur.

% =========================================================
\section{Numerical Examples}

\subsection{Numerical Example 1: Convolution}
Consider the input matrix:
\[
X = \begin{bmatrix} 1 & 2 & 3 \\ 4 & 5 & 6 \\ 7 & 8 & 9 \end{bmatrix}, \qquad
K = \begin{bmatrix} 1 & 0 \\ 0 & 1 \end{bmatrix}
\]

For the first position: $1(1) + 2(0) + 4(0) + 5(1) = 6$

For the second position: $2(1) + 3(0) + 5(0) + 6(1) = 8$

For the third position: $4(1) + 5(0) + 7(0) + 8(1) = 12$

For the fourth position: $5(1) + 6(0) + 8(0) + 9(1) = 14$

Therefore, the feature map is:
\[
Y = \begin{bmatrix} 6 & 8 \\ 12 & 14 \end{bmatrix}
\]

\subsection{Numerical Example 2: Max Pooling}
Consider the feature map:
\[
X = \begin{bmatrix} 1 & 5 & 2 & 3 \\ 7 & 8 & 1 & 0 \\ 4 & 6 & 9 & 5 \\ 2 & 3 & 1 & 8 \end{bmatrix}
\]
Using a $2 \times 2$ max pooling filter with stride 2:
\[
\begin{bmatrix} 1 & 5 \\ 7 & 8 \end{bmatrix} \rightarrow 8, \quad
\begin{bmatrix} 2 & 3 \\ 1 & 0 \end{bmatrix} \rightarrow 3, \quad
\begin{bmatrix} 4 & 6 \\ 2 & 3 \end{bmatrix} \rightarrow 6, \quad
\begin{bmatrix} 9 & 5 \\ 1 & 8 \end{bmatrix} \rightarrow 9
\]
Therefore:
\[
Y = \begin{bmatrix} 8 & 3 \\ 6 & 9 \end{bmatrix}
\]

\subsection{Numerical Example 3: Parameter Calculation}
Consider a convolution layer with:
\begin{itemize}
    \item Input size = $32 \times 32 \times 3$
    \item Number of filters = 16
    \item Kernel size = $3 \times 3$
\end{itemize}
The number of parameters in one filter is:
\[
3 \times 3 \times 3 + 1 = 28
\]
where the additional 1 represents the bias. Therefore:
\[
28 \times 16 = 448
\]
Thus, the convolution layer contains \textbf{448 trainable parameters}.

% =========================================================
\section{Dataset}
The CIFAR-10 dataset was used for this experiment. CIFAR-10 consists of 60,000 colour images divided into ten classes.

\begin{table}[h]
\centering
\caption{CIFAR-10 Dataset Characteristics}
\begin{tabular}{|l|l|}
\hline
\textbf{Property} & \textbf{Value} \\
\hline
Total images & 60,000 \\
Training images & 50,000 \\
Testing images & 10,000 \\
Number of classes & 10 \\
Image dimensions & $32 \times 32 \times 3$ \\
Image type & RGB \\
\hline
\end{tabular}
\end{table}

The ten classes are:
\begin{enumerate}
    \item Airplane
    \item Automobile
    \item Bird
    \item Cat
    \item Deer
    \item Dog
    \item Frog
    \item Horse
    \item Ship
    \item Truck
\end{enumerate}

The images were normalized by dividing their pixel values by 255:
\[
X_{\text{normalized}} = \frac{X}{255}
\]
This converts the pixel range from $[0, 255]$ to approximately $[0, 1]$.

% =========================================================
\section{Software and Tools}
\begin{table}[h]
\centering
\caption{Software Used}
\begin{tabular}{|l|l|}
\hline
\textbf{Tool/Library} & \textbf{Purpose} \\
\hline
Python & Programming language \\
Jupyter Notebook & Experimental environment \\
TensorFlow & Deep learning framework \\
Keras & Neural network API \\
NumPy & Numerical computation \\
Matplotlib & Data visualization \\
Seaborn & Confusion matrix visualization \\
Scikit-learn & Evaluation metrics \\
Pandas & Results organization \\
\hline
\end{tabular}
\end{table}

% =========================================================
\section{Experimental Procedure}
The experiment was performed using the following steps:
\begin{enumerate}
    \item Import the required Python libraries.
    \item Set random seeds for reproducibility.
    \item Load the CIFAR-10 dataset.
    \item Display dataset dimensions.
    \item Display ten sample images.
    \item Plot the class distribution.
    \item Normalize the input images.
    \item Perform a numerical convolution example.
    \item Compare $3\times3$, $5\times5$ and $7\times7$ convolution kernels.
    \item Study the effects of stride and padding.
    \item Implement common image-processing kernels.
    \item Implement max pooling and average pooling.
    \item Calculate convolutional layer parameters.
    \item Divide the training set into training and validation subsets.
    \item Construct the CNN architecture.
    \item Compile the model using the Adam optimizer.
    \item Train the model for 20 epochs with a batch size of 32.
    \item Plot training and validation accuracy.
    \item Plot training and validation loss.
    \item Visualize feature maps from the first convolutional layer.
    \item Evaluate the model on the test dataset.
    \item Calculate precision, recall and F1-score.
    \item Generate the confusion matrix and classification report.
    \item Replace max pooling with average pooling and compare results.
    \item Increase the number of filters and analyze the effect on performance.
\end{enumerate}

% =========================================================
\section{CNN Architecture}
The CNN architecture used in the main experiment is:
\[
\text{Input} \rightarrow \text{Conv2D} \rightarrow \text{ReLU} \rightarrow \text{MaxPool} \rightarrow \text{Conv2D} \rightarrow \text{ReLU} \rightarrow \text{MaxPool} \rightarrow \text{Flatten} \rightarrow \text{Dense} \rightarrow \text{Softmax}
\]

\begin{table}[h]
\centering
\caption{CNN Architecture}
\begin{tabular}{|l|l|l|l|}
\hline
\textbf{Layer} & \textbf{Configuration} & \textbf{Output} & \textbf{Parameters} \\
\hline
Input & $32 \times 32 \times 3$ & $32 \times 32 \times 3$ & 0 \\
Conv2D & 16 filters, $3\times3$ & $32 \times 32 \times 16$ & 448 \\
ReLU & -- & $32 \times 32 \times 16$ & 0 \\
MaxPool & $2\times2$ & $16 \times 16 \times 16$ & 0 \\
Conv2D & 32 filters, $3\times3$ & $16 \times 16 \times 32$ & 4,640 \\
ReLU & -- & $16 \times 16 \times 32$ & 0 \\
MaxPool & $2\times2$ & $8 \times 8 \times 32$ & 0 \\
Flatten & -- & 2048 & 0 \\
Dense & 128 neurons & 128 & 262,272 \\
ReLU & -- & 128 & 0 \\
Dense & 10 neurons & 10 & 1,290 \\
Softmax & -- & 10 & 0 \\
\hline
\textbf{Total} & & & \textbf{268,650} \\
\hline
\end{tabular}
\end{table}

The exact total parameter count can be obtained directly using \texttt{model.count\_params()} in the Jupyter notebook. The value above (268,650) is computed directly from the layer-by-layer arithmetic ($448 + 4{,}640 + 262{,}272 + 1{,}290$) and is exact, not an estimate.

% =========================================================
\section{Implementation}

\subsection{Importing Libraries}
\begin{lstlisting}
import os
import random
import numpy as np
import matplotlib.pyplot as plt
import seaborn as sns
import tensorflow as tf

from tensorflow import keras
from tensorflow.keras import layers

from sklearn.metrics import (
    accuracy_score,
    precision_score,
    recall_score,
    f1_score,
    confusion_matrix,
    classification_report
)
\end{lstlisting}

\subsection{Loading CIFAR-10}
\begin{lstlisting}
(x_train, y_train), (x_test, y_test) = keras.datasets.cifar10.load_data()

y_train = y_train.squeeze()
y_test = y_test.squeeze()

print(x_train.shape)
print(y_train.shape)
print(x_test.shape)
print(y_test.shape)
\end{lstlisting}

\subsection{Data Normalization}
\begin{lstlisting}
x_train_normalized = x_train.astype("float32") / 255.0
x_test_normalized = x_test.astype("float32") / 255.0
\end{lstlisting}

\subsection{Convolution Output Calculation}
\begin{lstlisting}
def conv_output_size(n, f, p, s):
    return ((n - f + 2 * p) // s) + 1

def same_output_size(n, s):
    return int(np.ceil(n / s))
\end{lstlisting}

\subsection{CNN Model}
\begin{lstlisting}
def build_cnn(pooling_type="max"):

    if pooling_type == "max":
        PoolingLayer = layers.MaxPooling2D
    elif pooling_type == "avg":
        PoolingLayer = layers.AveragePooling2D
    else:
        raise ValueError(
            "pooling_type must be 'max' or 'avg'"
        )

    model = keras.Sequential([
        layers.Input(shape=(32, 32, 3)),

        layers.Conv2D(
            filters=16,
            kernel_size=3,
            padding="same"
        ),

        layers.ReLU(),

        PoolingLayer(
            pool_size=2,
            strides=2
        ),

        layers.Conv2D(
            filters=32,
            kernel_size=3,
            padding="same"
        ),

        layers.ReLU(),

        PoolingLayer(
            pool_size=2,
            strides=2
        ),

        layers.Flatten(),

        layers.Dense(128),
        layers.ReLU(),

        layers.Dense(
            10,
            activation="softmax"
        )
    ])

    return model
\end{lstlisting}

\subsection{Model Compilation}
\begin{lstlisting}
model = build_cnn(pooling_type="max")

model.compile(
    optimizer=keras.optimizers.Adam(),
    loss="sparse_categorical_crossentropy",
    metrics=["accuracy"]
)
\end{lstlisting}

\subsection{Model Training}
\begin{lstlisting}
history = model.fit(
    x_training,
    y_training,
    validation_data=(x_validation, y_validation),
    epochs=20,
    batch_size=32,
    verbose=1
)
\end{lstlisting}

\subsection{Model Evaluation}
\begin{lstlisting}
test_loss, test_accuracy = model.evaluate(
    x_test_normalized,
    y_test,
    batch_size=32
)
\end{lstlisting}

\subsection{Prediction and Classification Metrics}
\begin{lstlisting}
probabilities = model.predict(
    x_test_normalized,
    batch_size=32
)

y_pred = np.argmax(probabilities, axis=1)

accuracy = accuracy_score(y_test, y_pred)

precision = precision_score(
    y_test, y_pred, average="weighted", zero_division=0
)

recall = recall_score(
    y_test, y_pred, average="weighted", zero_division=0
)

f1 = f1_score(
    y_test, y_pred, average="weighted", zero_division=0
)
\end{lstlisting}

% =========================================================
\section{Results}
The CNN was trained for 20 epochs using the Adam optimizer and a batch size of 32.

\begin{estbox}
The values in Table 4 were \textbf{not} produced by an actual notebook run. They are illustrative estimates, chosen to be internally consistent with the accuracy/loss curves and confusion matrix used later in this report, and typical of a two-conv-block CNN of this size on CIFAR-10 after 20 epochs without data augmentation or dropout (where mild overfitting -- training accuracy noticeably above validation/test accuracy -- is expected).
\end{estbox}

\begin{table}[h]
\centering
\caption{CNN Performance (illustrative estimate)}
\begin{tabular}{|l|l|}
\hline
\textbf{Metric} & \textbf{Value} \\
\hline
Training Accuracy & 0.812 \\
Validation Accuracy & 0.681 \\
Testing Accuracy & 0.676 \\
Precision (weighted) & 0.679 \\
Recall (weighted) & 0.676 \\
F1-score (weighted) & 0.676 \\
Number of Parameters & 268,650 \\
\hline
\end{tabular}
\end{table}

% =========================================================
\section{Mandatory Plots}

\subsection{Sample Images}
\begin{figure}[h]
\centering
\includegraphics[width=0.75\textwidth]{figs/sample_images.png}
\caption{Placeholder colour swatches standing in for ten sample CIFAR-10 images (illustrative only -- see note below).}
\end{figure}

\begin{estbox}
The image above is \textbf{not} drawn from the actual CIFAR-10 dataset -- it is a set of flat colour swatches used purely as a structural placeholder, since the real pixel data was not available when generating this document. Replace this figure with the actual \texttt{imshow} output from your notebook's Task/Step 5 before submission.
\end{estbox}

\textit{Inference:} The sample images demonstrate the diversity of objects present in the CIFAR-10 dataset. Each image has a spatial resolution of $32 \times 32$ pixels and contains three RGB channels.

\subsection{Class Distribution}
\begin{figure}[h]
\centering
\includegraphics[width=0.7\textwidth]{figs/class_distribution.png}
\caption{Class distribution of the CIFAR-10 training dataset.}
\end{figure}

\textit{Inference:} The CIFAR-10 training dataset is exactly balanced, with exactly 5,000 images in each of the ten classes ($50{,}000 / 10$). Therefore, significant class imbalance is not expected to bias the CNN towards a particular class. (This 5,000-per-class figure is an exact, well-known property of CIFAR-10, not an estimate.)

\subsection{Training Accuracy}
\begin{figure}[h]
\centering
\includegraphics[width=0.65\textwidth]{figs/training_accuracy.png}
\caption{Training accuracy over 20 epochs (illustrative estimate).}
\end{figure}

\textit{Inference:} Training accuracy generally increases as the CNN learns increasingly useful visual representations from the training images. A plateau indicates that additional training provides diminishing improvements on the training set.

\subsection{Validation Accuracy}
\begin{figure}[h]
\centering
\includegraphics[width=0.65\textwidth]{figs/validation_accuracy.png}
\caption{Validation accuracy over 20 epochs (illustrative estimate).}
\end{figure}

\textit{Inference:} Validation accuracy measures the generalization ability of the model on data not used for parameter updates. A large gap between training and validation accuracy can indicate overfitting -- as illustrated here, where validation accuracy plateaus well below training accuracy.

\subsection{Training Loss}
\begin{figure}[h]
\centering
\includegraphics[width=0.65\textwidth]{figs/training_loss.png}
\caption{Training loss over 20 epochs (illustrative estimate).}
\end{figure}

\textit{Inference:} Training loss generally decreases as the optimizer updates the model parameters to minimize the classification error. A decreasing loss indicates that the model is fitting the training data more effectively.

\subsection{Validation Loss}
\begin{figure}[h]
\centering
\includegraphics[width=0.65\textwidth]{figs/validation_loss.png}
\caption{Validation loss over 20 epochs (illustrative estimate).}
\end{figure}

\textit{Inference:} Validation loss provides an indication of how well the model generalizes to unseen data. In this illustrative curve, validation loss bottoms out early and then increases while training loss continues decreasing, which is a sign of overfitting.

\subsection{Feature Maps}
The output of the first convolutional layer was used to visualize the activation of eight learned filters.

\begin{figure}[h]
\centering
\includegraphics[width=0.75\textwidth]{figs/feature_maps.png}
\caption{Eight feature maps generated by the first convolutional layer (illustrative synthetic patterns).}
\end{figure}

\begin{estbox}
These feature maps are synthetic mathematical patterns (sinusoids, gradients, radial functions) generated purely to illustrate what a feature-map grid layout looks like. They are \textbf{not} real activations from a trained convolutional filter. Replace with the actual output of \texttt{feature\_extractor.predict(sample)} from Section 13 before submission.
\end{estbox}

\textit{Inference:} Each feature map represents the response of a learned convolutional filter to different visual patterns in the input image. Early convolutional layers typically learn low-level features such as edges, textures and colour transitions.

\subsection{Confusion Matrix}
\begin{figure}[h]
\centering
\includegraphics[width=0.75\textwidth]{figs/confusion_matrix.png}
\caption{Confusion matrix for CIFAR-10 test predictions (illustrative estimate, internally consistent with the 0.676 testing accuracy reported in Table 4).}
\end{figure}

\textit{Inference:} The confusion matrix shows correct predictions along the diagonal and misclassifications outside the diagonal. In this illustrative matrix, the largest off-diagonal confusions are placed between Cat/Dog and Automobile/Truck, consistent with well-documented, commonly-observed CIFAR-10 confusion patterns for small CNNs, since these class pairs share similar shapes, textures and colour distributions at $32\times32$ resolution.

% =========================================================
\section{Feature Map Visualization}
The first convolutional layer of the trained CNN was used to extract feature maps. The feature extractor was created using:
\begin{lstlisting}
feature_extractor = keras.Model(
    inputs=model.input,
    outputs=model.layers[0].output
)

sample = x_test_normalized[0:1]

feature_maps = feature_extractor.predict(
    sample,
    verbose=0
)
\end{lstlisting}

Eight feature maps were then visualized to inspect the patterns detected by the first convolutional layer.

The visualization demonstrates that convolutional filters transform the original image into different representations. Each representation highlights different aspects of the image.

% =========================================================
\section{Max Pooling vs Average Pooling}
To study the effect of pooling strategy, a second CNN was constructed by replacing max pooling with average pooling.

Both models used the same:
\begin{itemize}
    \item Dataset
    \item Number of convolutional layers
    \item Number of filters
    \item Kernel sizes
    \item Dense layers
    \item Optimizer
    \item Number of epochs
    \item Batch size
\end{itemize}
Only the pooling operation was changed.

\begin{estbox}
Both rows of Table 5 have identical parameter counts because pooling layers contain no trainable weights -- this is an exact, architecture-derived fact, not an estimate. The test-accuracy values themselves are illustrative estimates: max pooling is given a slight edge, consistent with the commonly-reported tendency of max pooling to preserve the sharp, high-contrast edge information that matters most for small natural-image datasets like CIFAR-10.
\end{estbox}

\begin{table}[h]
\centering
\caption{Pooling Comparison (illustrative estimate)}
\begin{tabular}{|l|l|l|}
\hline
\textbf{Pooling Method} & \textbf{Test Accuracy} & \textbf{Parameters} \\
\hline
Max Pooling & 0.676 & 268,650 \\
Average Pooling & 0.658 & 268,650 \\
\hline
\end{tabular}
\end{table}

\textit{Inference:} Max pooling retains the strongest activation within each local region, thereby emphasizing the most prominent features. Average pooling instead aggregates all values within the region and produces a smoother representation.

The better-performing pooling strategy is determined experimentally based on the resulting test accuracy; in this illustrative comparison, max pooling outperforms average pooling by roughly 1.8 percentage points.

% =========================================================
\section{Effect of Number of Filters}
The effect of increasing the number of convolutional filters was studied by comparing models using 16 and 64 filters (with the second convolutional layer scaled proportionally, from 32 to 128 filters, to keep the two-block architecture consistent).

Increasing the number of filters increases the representational capacity of the convolutional layers because more distinct visual patterns can be learned. However, it also increases the number of trainable parameters and therefore increases computational requirements.

\begin{estbox}
Parameters for the 64-filter model are computed exactly from the architecture: Conv2D(64): $(3\times3\times3+1)\times64 = 1{,}792$; Conv2D(128): $(3\times3\times64+1)\times128 = 73{,}856$; Flatten $8\times8\times128 = 8{,}192$; Dense(128): $8{,}192\times128+128 = 1{,}048{,}704$; Dense(10): $128\times10+10=1{,}290$. Total $= 1{,}125{,}642$ (exact). The test-accuracy values are illustrative estimates only, reflecting the typical (but not guaranteed) modest accuracy gain from a $4\times$ increase in filter count on a small 32$\times$32 dataset before diminishing returns and overfitting set in.
\end{estbox}

\begin{table}[h]
\centering
\caption{Effect of Number of Filters (illustrative estimate for accuracy; parameter counts exact)}
\begin{tabular}{|l|l|l|}
\hline
\textbf{Model} & \textbf{Test Accuracy} & \textbf{Parameters} \\
\hline
16 Filters (base: 16 / 32) & 0.676 & 268,650 \\
64 Filters (scaled: 64 / 128) & 0.712 & 1,125,642 \\
\hline
\end{tabular}
\end{table}

\textit{Inference:} Increasing the number of filters provides the network with additional capacity to learn visual representations. However, more filters result in higher memory usage, computation time and parameter count, and do not necessarily guarantee proportional improvements in classification accuracy -- here, a roughly $4.2\times$ increase in parameters yields an illustrative gain of only about 3.6 percentage points in accuracy.

% =========================================================
\section{Additional Exercises}

\subsection{Output Size for a 64 $\times$ 64 Image}
Given: $N = 64$, $F = 5$, $S = 2$, $P = 2$

Using:
\[
O = \left\lfloor \frac{N - F + 2P}{S} \right\rfloor + 1
\]
we obtain:
\[
O = \left\lfloor \frac{64 - 5 + 2(2)}{2} \right\rfloor + 1 = \left\lfloor \frac{63}{2} \right\rfloor + 1 = 31 + 1 = 32
\]
Therefore: $32 \times 32$.

\subsection{Parameters of a 64-Filter Convolution Layer}
Given:
\begin{itemize}
    \item Number of filters = 64
    \item Kernel = $3 \times 3$
    \item Input channels = 3
\end{itemize}
The number of parameters is:
\[
(3 \times 3 \times 3 + 1) \times 64 = (27 + 1) \times 64 = 28 \times 64 = 1792
\]
Therefore, the convolutional layer contains \textbf{1792 trainable parameters}.

\subsection{ReLU vs Sigmoid}
\begin{table}[h]
\centering
\caption{Comparison of ReLU and Sigmoid}
\begin{tabular}{|l|l|l|}
\hline
\textbf{Property} & \textbf{ReLU} & \textbf{Sigmoid} \\
\hline
Function & $\max(0, x)$ & $\frac{1}{1+e^{-x}}$ \\
Output range & $[0, \infty)$ & $(0, 1)$ \\
Computation & Simple & More expensive \\
Vanishing gradient & Lower risk & Higher risk \\
Typical CNN usage & Hidden layers & Binary output layers \\
\hline
\end{tabular}
\end{table}

ReLU is generally preferred in hidden CNN layers because of its computational simplicity and improved gradient behaviour. Sigmoid is more appropriate when an output probability between zero and one is required, particularly for binary classification.

% =========================================================
\section{Discussion}

\subsection{Why is convolution preferred over fully connected layers for images?}
Convolution is preferred because it exploits the spatial structure of images. A convolutional filter operates on local regions rather than connecting every input pixel to every neuron.

Additionally, convolution uses shared weights, meaning the same filter is applied at different spatial positions.

Consequently, CNNs require significantly fewer parameters than fully connected networks while still being able to learn meaningful spatial features.

\subsection{How does stride affect the feature map size?}
Stride controls the distance by which the convolution kernel moves across the input. A larger stride results in fewer kernel positions and therefore a smaller feature map.

For example, with a $3 \times 3$ kernel and valid padding:
\[
32 \times 32 \xrightarrow{S=1} 30 \times 30
\]
whereas:
\[
32 \times 32 \xrightarrow{S=2} 15 \times 15
\]
Thus, increasing stride reduces the spatial resolution of the feature map.

\subsection{What is the role of padding?}
Padding adds pixels around the boundaries of the input. Its main purposes are:
\begin{itemize}
    \item Preserve spatial dimensions.
    \item Allow boundary pixels to participate in convolution more frequently.
    \item Control the spatial size of feature maps.
    \item Prevent excessive reduction in image dimensions across multiple convolutional layers.
\end{itemize}

\subsection{Why is pooling used?}
Pooling reduces the spatial dimensions of feature maps. This results in:
\begin{itemize}
    \item Lower computational cost.
    \item Reduced memory requirements.
    \item More compact feature representations.
    \item Increased robustness to small spatial translations.
\end{itemize}
Max pooling retains the strongest activation, while average pooling computes the average activation within a local region.

\subsection{How do feature maps represent image characteristics?}
A feature map represents the response of a convolutional filter to an input image. Different filters learn different visual characteristics. Early CNN layers typically detect low-level features such as:
\begin{itemize}
    \item Edges
    \item Corners
    \item Colour transitions
    \item Simple textures
\end{itemize}
As information moves through deeper layers, the network combines these features to represent increasingly complex structures.

\subsection{Why do CNNs require fewer parameters than MLPs?}
CNNs use two important concepts:
\begin{enumerate}
    \item Local connectivity
    \item Weight sharing
\end{enumerate}
In an MLP, each neuron can be connected to every input pixel. For a large image, this produces a very large number of parameters.

In contrast, a CNN applies a relatively small filter repeatedly across the image. For example, a $3 \times 3$ filter operating on an RGB image requires only:
\[
3 \times 3 \times 3 = 27
\]
weights plus one bias. The same 28 parameters are reused across the entire spatial image. This dramatically reduces the number of trainable parameters.

% =========================================================
\section{Overall Observations}
The following observations were made during the experiment:
\begin{enumerate}
    \item Kernel size determines the spatial region considered by each convolutional operation.
    \item Larger kernels capture larger receptive fields but increase computational requirements.
    \item Increasing stride reduces feature-map spatial dimensions.
    \item Same padding helps preserve spatial dimensions.
    \item Max pooling preserves the strongest activation in a local region.
    \item Average pooling produces smoother representations by averaging activations.
    \item Feature maps provide an interpretable representation of the intermediate features learned by the CNN.
    \item Increasing the number of filters increases the representational capacity of a CNN.
    \item Increasing the number of filters also increases parameter count, memory usage and computational requirements.
    \item Training and validation curves can be used to identify possible overfitting.
    \item The confusion matrix provides class-level information about classification performance.
\end{enumerate}

% =========================================================
\section{Final Result}
A Convolutional Neural Network was successfully implemented using TensorFlow/Keras for CIFAR-10 image classification.

The experiment successfully demonstrated:
\begin{itemize}
    \item Convolution operations.
    \item Convolution kernel behaviour.
    \item Output dimension calculation.
    \item Effects of stride and padding.
    \item Max and average pooling.
    \item CNN parameter calculation.
    \item Feature-map visualization.
    \item CNN model training.
    \item Image classification.
    \item Performance evaluation.
    \item Pooling strategy comparison.
    \item Filter-count comparison.
\end{itemize}

The final numerical results obtained from the experiment are shown below.

\begin{table}[h]
\centering
\caption{Final Experimental Results (illustrative estimate -- see note in Section 11)}
\begin{tabular}{|l|l|}
\hline
\textbf{Metric} & \textbf{Result} \\
\hline
Training Accuracy & 0.812 \\
Validation Accuracy & 0.681 \\
Testing Accuracy & 0.676 \\
Precision & 0.679 \\
Recall & 0.676 \\
F1-score & 0.676 \\
Total Parameters & 268,650 \\
\hline
\end{tabular}
\end{table}

% =========================================================
\section{Conclusion}
In this experiment, a Convolutional Neural Network was implemented using TensorFlow/Keras for image classification on the CIFAR-10 dataset.

The experiment provided practical understanding of convolution, activation, pooling and fully connected layers. The effect of kernel size, stride and padding on feature-map dimensions was experimentally studied. Classical convolution kernels were also implemented to demonstrate how filters can extract edges, textures, boundaries and other image characteristics.

A CNN consisting of two convolutional layers, ReLU activation, pooling, flattening and dense classification layers was trained using the Adam optimizer for 20 epochs with a batch size of 32.

Feature-map visualization demonstrated how convolutional filters transform input images into intermediate representations. Max pooling and average pooling were compared experimentally, and the effect of increasing the number of convolutional filters was also investigated.

The CNN was evaluated using accuracy, precision, recall, F1-score, classification report and confusion matrix. Based on the illustrative estimates used in this report (Section 11), the baseline configuration reached approximately 67.6\% test accuracy with 268,650 trainable parameters, with max pooling and a larger 64/128-filter configuration both showing modest expected improvements. Overall, the experiment demonstrated why CNNs are particularly effective for image classification: they exploit spatial locality and weight sharing while progressively learning hierarchical visual representations.

\begin{estbox}
As stated throughout this document, the specific accuracy/precision/recall/F1 values, the confusion matrix, the training curves, the sample images and the feature maps are illustrative estimates, not real notebook output. Re-run the code in Appendix A and replace the flagged figures/tables with your actual results before submission.
\end{estbox}

% =========================================================
\section{References}
\begin{enumerate}
    \item Ian Goodfellow, Yoshua Bengio and Aaron Courville, \textit{Deep Learning}, MIT Press, 2016.
    \item Christopher M. Bishop, \textit{Pattern Recognition and Machine Learning}, Springer, 2006.
    \item Simon Haykin, \textit{Neural Networks and Learning Machines}, Pearson.
    \item TensorFlow Documentation, TensorFlow and Keras API documentation.
    \item CIFAR-10 Dataset Documentation, Canadian Institute for Advanced Research.
\end{enumerate}

% =========================================================
\appendix
\section{Complete CNN Implementation}
The following code represents the primary CNN implementation used in the experiment.

\begin{lstlisting}
import os
import random
import numpy as np
import matplotlib.pyplot as plt
import seaborn as sns
import tensorflow as tf

from tensorflow import keras
from tensorflow.keras import layers

from sklearn.metrics import (
    accuracy_score,
    precision_score,
    recall_score,
    f1_score,
    confusion_matrix,
    classification_report
)

# Reproducibility
SEED = 42

os.environ["PYTHONHASHSEED"] = str(SEED)

random.seed(SEED)
np.random.seed(SEED)
tf.random.set_seed(SEED)

# Load dataset
(x_train, y_train), (x_test, y_test) = \
    keras.datasets.cifar10.load_data()

y_train = y_train.squeeze()
y_test = y_test.squeeze()

# Normalize
x_train_normalized = \
    x_train.astype("float32") / 255.0

x_test_normalized = \
    x_test.astype("float32") / 255.0

# Validation split
VALIDATION_SIZE = 5000

x_validation = \
    x_train_normalized[-VALIDATION_SIZE:]

y_validation = \
    y_train[-VALIDATION_SIZE:]

x_training = \
    x_train_normalized[:-VALIDATION_SIZE]

y_training = \
    y_train[:-VALIDATION_SIZE]

# Build model
model = keras.Sequential([

    layers.Input(shape=(32, 32, 3)),

    layers.Conv2D(
        filters=16,
        kernel_size=3,
        padding="same"
    ),

    layers.ReLU(),

    layers.MaxPooling2D(
        pool_size=2,
        strides=2
    ),

    layers.Conv2D(
        filters=32,
        kernel_size=3,
        padding="same"
    ),

    layers.ReLU(),

    layers.MaxPooling2D(
        pool_size=2,
        strides=2
    ),

    layers.Flatten(),

    layers.Dense(128),

    layers.ReLU(),

    layers.Dense(
        10,
        activation="softmax"
    )
])

# Compile
model.compile(
    optimizer=keras.optimizers.Adam(),
    loss="sparse_categorical_crossentropy",
    metrics=["accuracy"]
)

# Train
history = model.fit(
    x_training,
    y_training,
    validation_data=(
        x_validation,
        y_validation
    ),
    epochs=20,
    batch_size=32,
    verbose=1
)

# Evaluate
test_loss, test_accuracy = model.evaluate(
    x_test_normalized,
    y_test,
    batch_size=32
)

# Predictions
probabilities = model.predict(
    x_test_normalized,
    batch_size=32
)

y_pred = np.argmax(
    probabilities,
    axis=1
)

# Metrics
accuracy = accuracy_score(
    y_test,
    y_pred
)

precision = precision_score(
    y_test,
    y_pred,
    average="weighted",
    zero_division=0
)

recall = recall_score(
    y_test,
    y_pred,
    average="weighted",
    zero_division=0
)

f1 = f1_score(
    y_test,
    y_pred,
    average="weighted",
    zero_division=0
)

print("Accuracy:", accuracy)
print("Precision:", precision)
print("Recall:", recall)
print("F1-score:", f1)

# Classification report
print(
    classification_report(
        y_test,
        y_pred,
        digits=4
    )
)

# Confusion matrix
cm = confusion_matrix(
    y_test,
    y_pred
)

plt.figure(figsize=(11, 9))

sns.heatmap(
    cm,
    annot=True,
    fmt="d",
    cmap="Blues"
)

plt.title(
    "CIFAR-10 Confusion Matrix"
)

plt.xlabel(
    "Predicted Class"
)

plt.ylabel(
    "True Class"
)

plt.tight_layout()
plt.show()
\end{lstlisting}

\end{document}
